\documentclass[../../../include/open-logic-section]{subfiles}

\begin{document}

\olfileid[hi]{sth}{replacement}{extrinsic}
\olsection[बाह्य विचार]{प्रतिस्थापन के विषय में बाह्य विचार}

हम प्रतिस्थापन का औचित्य देने के एक \emph{बाह्य} प्रयास पर विचार करने से
आरंभ करते हैं। बूलोस निम्न सुझाव देते हैं।
\begin{quote}
  [\ldots] प्रतिस्थापन के अभिगृहीत अपनाने का कारण काफ़ी सरल है: उनके अनेक
  वांछनीय परिणाम हैं और (प्रकटतः) कोई अवांछनीय परिणाम नहीं। पुनरावर्ती
  अवधारणा से संबंधित प्रमेयों के अतिरिक्त परिणामों में अनंत संख्याओं का
  संतोषजनक, भले ही आदर्श न हो, सिद्धांत और सुस्थापित संबंधों पर आगमनात्मक
  परिभाषाओं का औचित्य देने वाला अत्यंत वांछनीय परिणाम सम्मिलित हैं।
  \citep[229]{Boolos1971}
\end{quote}		
बूलोस के विचार का सार है कि हमें प्रतिस्थापन का औचित्य उसके फलों से देना
चाहिए। उनके बताए विशिष्ट फल वे बातें हैं जिनकी चर्चा हमने पिछले कुछ अध्यायों
में की है। प्रतिस्थापन ने हमें सिद्ध करने दिया कि वॉन नोयमान क्रमसूचक
सुव्यवस्था-प्रकार के विचार के उत्कृष्ट प्रतिनिधि थे (यही हमारा ``अनंत
संख्याओं का संतोषजनक, भले ही आदर्श न हो, सिद्धांत'' है)। प्रतिस्थापन ने हमें
$V_\alpha$ परिभाषित करने, रैंक की धारणा स्थापित करने और $\in$-आगमन सिद्ध
करने दिया (यही ``पुनरावर्ती अवधारणा के विषय में हमारे प्रमेय'' हैं)। अंततः
प्रतिस्थापन हमें परिमितेतर पुनरावर्तन प्रमेय सिद्ध करने देता है (यही
``सुस्थापित संबंधों पर आगमनात्मक परिभाषाएँ'' हैं)।

ये वास्तव में वांछनीय परिणाम हैं। किंतु क्या ये वांछनीय परिणाम प्रतिस्थापन
का \emph{औचित्य} देने के लिए पर्याप्त हैं? \emph{नहीं।} कम-से-कम सीधे नहीं।

एक सरल समस्या यह है। हमने प्रतिस्थापन के कुछ वांछनीय परिणाम बताए हैं, किंतु
उनमें से कई अन्य साधनों से प्राप्त किए जा सकते थे। यह बात जितनी प्रसिद्ध
होनी चाहिए उतनी नहीं, इसलिए स्थिति समझाने के लिए रुकना उचित है।

समुच्चयों का एक सरल सिद्धांत, स्तर सिद्धांत या संक्षेप में $\LT$, है।
\footnote{$\LT$ के पहले रूप \citet{Montague1965} और \citet{Scott1974} ने
दिए; \citet{Potter2004} ने उसे सरल बनाकर पुस्तक-लंबा विवेचन दिया; और
\citet{ButtonLT1} ने हाल में $\LT$ को और सरल बनाया है।} $\LT$ के अभिगृहीत
केवल विस्तारिता, पृथक्करण और यह दावा हैं कि प्रत्येक समुच्चय किसी
\emph{स्तर} का उपसमुच्चय है, जहाँ ``स्तर'' को चतुराई से इस प्रकार परिभाषित
किया जाता है कि स्तर हमारे परिचित $V_\alpha$ जैसा व्यवहार करें। अतः $\ZF$,
$\LT$ को सिद्ध करता है; किंतु $\LT$, $\ZF$ से \emph{बहुत} दुर्बल है। वास्तव
में $\LT$ युग्म, घात समुच्चय, अनंतता या प्रतिस्थापन नहीं देता। $\LT$ में
अनंतता और घात समुच्चय जोड़ने से प्राप्त सिद्धांत को $\Zr$ मानें; इससे युग्म
भी मिलता है, अतः $\Zr$ कम-से-कम $\Z$ जितना प्रबल है। किंतु वास्तव में
$\Zr$, $\Z$ से कड़ाई से अधिक प्रबल है, क्योंकि वह यह दावा जोड़ता है कि
प्रत्येक समुच्चय की कोई रैंक है (इसीलिए इसे $\Zr$ कहने का सुझाव)। वास्तव में
$\Zr$ देता है: क्रमसूचकों का पूर्णतः संतोषजनक सिद्धांत; पदानुक्रम को
सुव्यवस्थित चरणों में स्तरित करने वाले परिणाम; $\in$-आगमन का प्रमाण; और
परिमितेतर पुनरावर्तन का एक \emph{रूप}।

संक्षेप में: यद्यपि बूलोस यह नहीं जानते थे, उनके बताए सभी वांछनीय परिणाम
प्रतिस्थापन के \emph{बिना} प्राप्त किए जा सकते थे; उन्हें बस $\Z$ के बजाय
$\Zr$ का उपयोग करना था।

(यह सब देखते हुए हमने आपको $\LT$ और $\Zr$ के बजाय $\ZF$ पढ़ाने का पारंपरिक
मार्ग क्यों अपनाया? दो कारण हैं। पहला: विशुद्ध ऐतिहासिक कारणों से $\LT$ से
आरंभ करना काफ़ी अमानक है; हम आपको समुच्चय सिद्धांत की अधिक मानक चर्चाएँ पढ़
सकने में समर्थ बनाना चाहते थे। दूसरा: जब आप $\LT$ और $\Zr$ को समझने के लिए
तैयार हों, तब बस \citealt{Potter2004} और \citealt{ButtonLT1} पढ़ सकते हैं।)

निश्चय ही चूँकि $\Zr$, $\ZF$ से कड़ाई से दुर्बल है, कुछ परिणाम ऐसे हैं
जिन्हें $\ZF$ सिद्ध करता है किंतु $\Zr$ अनिर्णीत छोड़ता है। अतः इनमें से
किसी परिणाम की ओर संकेत करके बाह्य आधार पर प्रतिस्थापन का औचित्य देने का
प्रयास किया जा सकता है। किंतु $\Zr$ का उपयोग जान लेने पर ऐसी बातों के बहुत
उदाहरण खोजना कठिन है जो (a) प्रतिस्थापन द्वारा, किंतु अन्यथा नहीं, निर्णीत
होती हैं और (b) सहजतः सत्य हैं। (अधिक के लिए
\citealt[\S13.2]{Potter2004} देखें।)

निष्कर्ष यह है। प्रतिस्थापन का विश्वसनीय बाह्य औचित्य देने के लिए ऐसा परिणाम
खोजना होगा जो प्रतिस्थापन के बिना प्राप्त \emph{नहीं किया जा सकता}। और यह
सरल काम नहीं।

प्रतिस्थापन का विशुद्ध बाह्य औचित्य देने के प्रत्येक प्रयास में उत्पन्न एक
और समस्या पर विचार करें। (यह समस्या शायद पहली से अधिक मूलभूत है।) बूलोस
केवल यह नहीं बताते कि प्रतिस्थापन के अनेक वांछनीय परिणाम हैं। वे यह भी कहते
हैं कि उसके ``(प्रकटतः) कोई अवांछनीय'' परिणाम नहीं। किंतु कोष्ठक में दिया यह
संरक्षण, ``प्रकटतः'', निश्चय ही अत्यंत निर्णायक है।

याद करें कि हम यहाँ कैसे पहुँचे: सहज समुच्चय-निर्माण असंगति में पड़ा और हमने
इस असंगति का उत्तर समुच्चय की संचयी-पुनरावर्ती अवधारणा अपनाकर दिया। इस
अवधारणा के साथ ऐसी कथा आती है जो, हमें आशा है, उसकी संगति का आश्वासन देती
है। किंतु यदि उस कथा के भीतर से प्रतिस्थापन का औचित्य न दिया जा सके, तो
$\ZF$ को संगत मानने का (अभी तक) कोई कारण नहीं। अथवा अधिक ठीक कहें तो शायद
केवल संयोगवश अब तक किसी ने कोई विरोधाभास \emph{नहीं खोजा}, इस तथ्य के
अतिरिक्त $\ZF$ को संगत मानने का कोई कारण नहीं। ठीक इसी अर्थ में बूलोस की
टिप्पणी का सार यह प्रतीत होता है: ``(प्रकटतः) $\ZF$ संगत है''। हमें संगति
का इससे अधिक आश्वासन माँगना चाहिए।

यह समस्या प्रतिस्थापन का औचित्य देने के प्रत्येक \emph{विशुद्ध} बाह्य
प्रयास को प्रभावित करेगी, अर्थात ऐसे हर औचित्य को जो केवल $\ZF$ के (ज्ञात)
परिणामों के पदों में व्यक्त हो। इसलिए हम प्रतिस्थापन का \emph{आंतरिक}
औचित्य खोजना चाहेंगे, अर्थात ऐसा औचित्य जो सुझाए कि समुच्चयों के विषय में
हमारी कथा किसी तरह हमें प्रतिस्थापन के प्रति ``पहले से'' प्रतिबद्ध करती है।

\end{document}
